\section{Related Work}
\label{sec:related}



\niparagraph{LLM inference serving.}
%
Multiple LLM serving systems optimize for their inference performance either through reducing the memory footprint~\cite{arxiv21-lightseq, kvcaching}, improving the kernel execution strategy~\cite{alpa}, determining the partitioning techniques for intra- and inter- operator~\cite{pip, megatronlm, arxiv22-efficiently} execution, or a combination of these~\cite{deepspeed, orca, vllm, hugging-face, tensorflow-serving, tensorrt-llm, nvidia_triton, miao2023spotserve, li2023alpaserve}.
%
In this work, we specifically tackle the utilization of the current hardware platforms which deploy these models through a compute and I/O suitable platforms, NPUs and PIMs, to create a more efficient system. 
%
Moreover, to ensure such a heterogeneous system can perform well for LLM inferencing, we offer a scheduling policy.
%
Prior works that support kernel optimizations for better utilization of GPUs for LLMs, cannot fully mitigate the I/O and bandwidth bottleneck of GEMV kernels. 
%
Instead in this work, we build a system that can benefit from existing optimizations such as selective batching, KV caching, etc. to offer better utilization across the transformer model architecture.
%

\niparagraph{PIM for language model support.}
%
TransPIM~\cite{hpca22-transpim} is a PIM solution that accelerates the end-to-end transformer inference using PIM. 
%
The work proposes a data loading overhead reduction technique by customizing its dataflow for transformer models.
%
TransPIM is optimized for attention operations of transformer encoder blocks, making it unsuitable for LLM inference based on decoder blocks.
%
Furthermore, it is tailored for single-request inference, offering suboptimal performance for batched inference scenarios.
%
AttAcc~\cite{AttAcc} further offers an accelerator (with PIM) for attention layer to reduce the data movement for the KV matrices.
%
Instead \neupims proposes a new system for PIM accelerator in addition to the scheduling of operations for the end to end inference of LLMs. 

There are also variety of prior works that leverage PIM for GEMV operations~\cite{hcs23-aquabolt-xl, isscc22-aim-gddr6, micro21-trim, hbm_pim, newton, tensor-dimm, SpaceA} due to their inherent potential in benefits towards bandwidth bound applications.
%
However, none of these works enable simultaneous execution of PIM and NPU operations, necessary for the efficient execution of LLM inference. 

\niparagraph{Heterogeneous acceleration pipeline for deep learning.}
%
There are a variety of prior works that propose a pipelined solution for machine learning~\cite{fae, dnnweaver:micro, cosmic:micro, dana}, however none of these prior works leverage PIM to alleviate the bandwidth requirement of LLMs.
%
Certain prior works use an accelerator for specific models~\cite{hotline, scratchpipe, tabla:hpca}, however, they do not alleviate the under-utilization of GEMV and GEMM operations in transformer decoder blocks.

%

%
\if 0
Samsung Aquabolt-XL\cite{hcs23-aquabolt-xl}
SK Hynix Aim-GDDR6~\cite{isscc22-aim-gddr6}
trim\cite{micro21-trim}
HBM-PIM \cite{hbm_pim}
Newton \cite{newton}
tensor-dimm \cite{tensor-dimm}
SpaceA\cite{SpaceA}
%
\fi 